\chapter{50 Pertanyaan \& Jawaban Wawancara AI}
\label{chap:interview}

Persiapan adalah kunci kesuksesan dalam wawancara kerja di bidang AI dan Data Science. Berikut adalah kompilasi 50 pertanyaan yang paling sering muncul, mulai dari konsep dasar hingga desain sistem tingkat lanjut.

\section{Machine Learning Fundamentals}

\subsection{1. Apa perbedaan antara Supervised, Unsupervised, dan Reinforcement Learning?}
\textbf{Jawaban:}
\begin{itemize}
    \item \textbf{Supervised Learning:} Model belajar dari data berlabel (input-output pairs). Tujuannya adalah memprediksi output untuk data baru. Contoh: Klasifikasi email spam, prediksi harga rumah (regresi).
    \item \textbf{Unsupervised Learning:} Model belajar dari data yang tidak berlabel. Tujuannya adalah menemukan pola atau struktur tersembunyi. Contoh: Segmentasi pelanggan (clustering), reduksi dimensi (PCA).
    \item \textbf{Reinforcement Learning:} Agen belajar berinteraksi dengan lingkungan melalui trial-and-error untuk memaksimalkan reward kumulatif. Contoh: AlphaGo, robotika, game AI.
\end{itemize}

\subsection{2. Jelaskan Bias-Variance Tradeoff!}
\textbf{Jawaban:}
Bias-Variance Tradeoff adalah masalah mendasar dalam supervised learning yang berkaitan dengan kemampuan generalisasi model.
\begin{itemize}
    \item \textbf{High Bias (Underfitting):} Model terlalu sederhana untuk menangkap pola dalam data. Error tinggi pada data latih dan data uji.
    \item \textbf{High Variance (Overfitting):} Model terlalu kompleks dan menghafal noise dalam data latih. Error rendah pada data latih, tapi tinggi pada data uji.
    \item \textbf{Tradeoff:} Kita ingin menyeimbangkan keduanya untuk meminimalkan total error. Solusinya bisa dengan regularisasi, menambah data, atau mengubah kompleksitas model.
\end{itemize}

\subsection{3. Apa itu Overfitting dan bagaimana cara mencegahnya?}
\textbf{Jawaban:}
Overfitting terjadi ketika model belajar detail dan noise pada data latih "terlalu baik" sehingga gagal menggeneralisasi ke data baru.
Cara mencegahnya:
\begin{itemize}
    \item Menambah lebih banyak data latih.
    \item Menggunakan teknik Regularisasi (L1/Lasso, L2/Ridge).
    \item Menggunakan Dropout (pada Neural Networks).
    \item Early Stopping (menghentikan pelatihan saat validation loss mulai naik).
    \item Mengurangi kompleksitas model (misal: mengurangi kedalaman pohon keputusan).
    \item Data Augmentation.
\end{itemize}

\subsection{4. Jelaskan perbedaan antara Precision dan Recall!}
\textbf{Jawaban:}
Keduanya adalah metrik evaluasi klasifikasi yang penting, terutama untuk dataset yang tidak seimbang (imbalanced).
\begin{itemize}
    \item \textbf{Precision:} Dari semua yang diprediksi Positif, berapa yang benar-benar Positif?
    \[ \text{Precision} = \frac{TP}{TP + FP} \]
    Penting ketika biaya False Positive tinggi (misal: filter spam, jangan sampai email penting masuk spam).
    \item \textbf{Recall (Sensitivity):} Dari semua yang benar-benar Positif, berapa yang berhasil ditemukan?
    \[ \text{Recall} = \frac{TP}{TP + FN} \]
    Penting ketika biaya False Negative tinggi (misal: deteksi kanker, jangan sampai ada pasien sakit yang dibilang sehat).
\end{itemize}

\subsection{5. Apa itu F1 Score?}
\textbf{Jawaban:}
F1 Score adalah rata-rata harmonik (harmonic mean) dari Precision dan Recall.
\[ F1 = 2 \times \frac{\text{Precision} \times \text{Recall}}{\text{Precision} + \text{Recall}} \]
Digunakan ketika kita ingin keseimbangan antara Precision dan Recall, dan dataset tidak seimbang.

\subsection{6. Apa itu Cross-Validation?}
\textbf{Jawaban:}
Teknik untuk menilai kinerja model dengan membagi data menjadi beberapa subset. Metode paling umum adalah K-Fold Cross-Validation:
1. Bagi data menjadi K bagian (folds).
2. Latih model pada K-1 bagian, uji pada 1 bagian sisanya.
3. Ulangi K kali hingga setiap bagian pernah menjadi data uji.
4. Rata-ratakan skor evaluasi.
Ini memberikan estimasi kinerja yang lebih robust daripada sekadar split train/test tunggal.

\subsection{7. Apa perbedaan antara Normalization dan Standardization?}
\textbf{Jawaban:}
\begin{itemize}
    \item \textbf{Normalization (Min-Max Scaling):} Mengubah nilai fitur ke rentang tertentu, biasanya [0, 1].
    \[ X_{new} = \frac{X - X_{min}}{X_{max} - X_{min}} \]
    Berguna ketika algoritma sensitif terhadap skala (misal: Neural Networks, KNN) dan kita tahu batas data.
    \item \textbf{Standardization (Z-Score Normalization):} Mengubah data sehingga memiliki rata-rata 0 dan standar deviasi 1.
    \[ X_{new} = \frac{X - \mu}{\sigma} \]
    Lebih tahan terhadap outlier dan disukai oleh algoritma seperti SVM dan Logistic Regression.
\end{itemize}

\subsection{8. Jelaskan cara kerja Decision Tree!}
\textbf{Jawaban:}
Decision Tree memecah data menjadi subset yang lebih kecil berdasarkan aturan keputusan if-then.
Prosesnya:
1. Pilih fitur terbaik untuk memisahkan data (menggunakan metrik seperti Gini Impurity atau Information Gain/Entropy).
2. Buat cabang berdasarkan nilai fitur tersebut.
3. Ulangi secara rekursif untuk setiap cabang hingga kriteria berhenti terpenuhi (misal: kedalaman maksimum tercapai, atau node murni).

\subsection{9. Apa itu Random Forest?}
\textbf{Jawaban:}
Random Forest adalah metode ensemble learning yang terdiri dari banyak Decision Trees.
Prinsipnya:
\begin{itemize}
    \item \textbf{Bagging (Bootstrap Aggregating):} Setiap pohon dilatih pada subset data acak (dengan pengembalian).
    \item \textbf{Feature Randomness:} Saat memecah node, hanya subset fitur acak yang dipertimbangkan.
\end{itemize}
Prediksi akhir diambil berdasarkan voting mayoritas (klasifikasi) atau rata-rata (regresi) dari semua pohon. Ini mengurangi varians dan overfitting dibandingkan decision tree tunggal.

\subsection{10. Apa perbedaan antara K-Means dan KNN?}
\textbf{Jawaban:}
\begin{itemize}
    \item \textbf{K-Means:} Algoritma \textit{Unsupervised} untuk \textit{Clustering}. Mengelompokkan data menjadi K cluster berdasarkan centroid. Membutuhkan proses pelatihan (mencari centroid).
    \item \textbf{KNN (K-Nearest Neighbors):} Algoritma \textit{Supervised} untuk \textit{Klasifikasi/Regresi}. Mengklasifikasikan data baru berdasarkan mayoritas label dari K tetangga terdekatnya. Tidak ada fase pelatihan eksplisit (lazy learner).
\end{itemize}

\section{Deep Learning \& Neural Networks}

\subsection{11. Apa itu Activation Function dan mengapa kita membutuhkannya?}
\textbf{Jawaban:}
Fungsi aktivasi menentukan output dari neuron. Tanpa fungsi aktivasi non-linear, Neural Network (seberapa pun dalamnya) hanya akan menjadi regresi linear raksasa. Non-linearitas memungkinkan network mempelajari pola kompleks.
Contoh: ReLU (umum untuk hidden layers), Sigmoid (output biner), Softmax (output multikelas).

\subsection{12. Jelaskan Vanishing Gradient Problem!}
\textbf{Jawaban:}
Terjadi pada deep network saat menggunakan fungsi aktivasi seperti Sigmoid atau Tanh. Selama backpropagation, gradien (turunan) dikalikan terus menerus dari output ke input layer. Karena turunan Sigmoid < 0.25, perkalian berulang menyebabkan gradien mendekati nol. Akibatnya, bobot di layer awal tidak pernah diperbarui, dan model gagal belajar. Solusi: Gunakan ReLU, LSTM, atau ResNet.

\subsection{13. Apa fungsi dari Optimizer (seperti Adam, SGD)?}
\textbf{Jawaban:}
Optimizer adalah algoritma yang memperbarui bobot model berdasarkan gradien loss function untuk meminimalkan error.
\begin{itemize}
    \item \textbf{SGD (Stochastic Gradient Descent):} Memperbarui bobot menggunakan satu sampel/batch kecil. Sederhana tapi bisa berosilasi.
    \item \textbf{Adam:} Mengadaptasi learning rate untuk setiap parameter menggunakan estimasi momen pertama (mean) dan kedua (variance) dari gradien. Standar de facto saat ini.
\end{itemize}

\subsection{14. Apa perbedaan antara Epoch, Batch, dan Iteration?}
\textbf{Jawaban:}
\begin{itemize}
    \item \textbf{Epoch:} Satu kali pass seluruh dataset melalui network.
    \item \textbf{Batch:} Sebagian kecil dataset yang diproses sekaligus.
    \item \textbf{Iteration:} Jumlah update bobot (jumlah batch) yang dibutuhkan untuk menyelesaikan satu epoch.
    Contoh: Dataset 1000 sampel, Batch Size 100. Maka 1 Epoch = 10 Iterasi.
\end{itemize}

\subsection{15. Jelaskan konsep Convolutional Neural Network (CNN)!}
\textbf{Jawaban:}
CNN dirancang untuk memproses data grid (gambar). Komponen utamanya:
1. \textbf{Convolutional Layer:} Menggunakan filter/kernel yang digeser ke seluruh gambar untuk mengekstrak fitur (tepi, tekstur).
2. \textbf{Pooling Layer (Max/Average):} Mengurangi dimensi spasial (downsampling) untuk mengurangi komputasi dan membuat fitur invarian terhadap pergeseran kecil.
3. \textbf{Fully Connected Layer:} Layer dense di akhir untuk klasifikasi.

\subsection{16. Apa itu Transfer Learning?}
\textbf{Jawaban:}
Teknik menggunakan model yang sudah dilatih pada dataset besar (misal: ImageNet) sebagai titik awal untuk tugas baru dengan data terbatas.
Kita bisa membekukan (freeze) layer awal (feature extraction) dan hanya melatih layer akhir (classifier), atau melakukan fine-tuning pada seluruh network dengan learning rate kecil.

\subsection{17. Jelaskan arsitektur Transformer!}
\textbf{Jawaban:}
Arsitektur deep learning yang diperkenalkan dalam paper "Attention Is All You Need" (2017). Mengandalkan mekanisme \textbf{Self-Attention} untuk memproses seluruh urutan input secara paralel (tidak sekuensial seperti RNN).
Komponen kunci: Multi-Head Attention, Positional Encoding, Feed-Forward Networks, dan Layer Normalization. Menjadi dasar bagi BERT, GPT, dan T5.

\subsection{18. Apa itu BERT vs GPT?}
\textbf{Jawaban:}
\begin{itemize}
    \item \textbf{BERT (Bidirectional Encoder Representations from Transformers):} Model \textit{Encoder-only}. Membaca teks dua arah (kiri-ke-kanan dan kanan-ke-kiri) sekaligus. Bagus untuk pemahaman (klasifikasi, QA, NER). Dilatih dengan Masked Language Modeling (MLM).
    \item \textbf{GPT (Generative Pre-trained Transformer):} Model \textit{Decoder-only}. Membaca teks satu arah (kiri-ke-kanan). Bagus untuk generasi teks. Dilatih dengan Causal Language Modeling (prediksi kata berikutnya).
\end{itemize}

\subsection{19. Apa itu Word Embedding (Word2Vec, GloVe)?}
\textbf{Jawaban:}
Representasi kata sebagai vektor dense berdimensi rendah (misal: 300-d) di mana kata-kata dengan makna semantik serupa berada berdekatan dalam ruang vektor. Memungkinkan operasi aljabar seperti: Raja - Pria + Wanita = Ratu.

\subsection{20. Jelaskan Long Short-Term Memory (LSTM)!}
\textbf{Jawaban:}
Jenis RNN khusus yang dirancang untuk mengatasi masalah vanishing gradient. Memiliki struktur "sel" dengan tiga gerbang (gates):
1. \textbf{Forget Gate:} Memutuskan informasi apa yang harus dibuang dari cell state.
2. \textbf{Input Gate:} Memutuskan informasi baru apa yang harus disimpan.
3. \textbf{Output Gate:} Menentukan output hidden state berdasarkan cell state.
Ini memungkinkan LSTM mengingat ketergantungan jangka panjang.

\section{LLM & Generative AI}

\subsection{21. Apa itu Tokenization dalam konteks LLM?}
\textbf{Jawaban:}
Proses memecah teks menjadi unit-unit kecil (token) yang bisa berupa kata, sebagian kata, atau karakter. LLM modern menggunakan \textit{Subword Tokenization} (seperti Byte-Pair Encoding atau WordPiece) untuk menangani kata-kata langka dan menjaga efisiensi. Contoh: "unbelievable" -> "un", "believ", "able".

\subsection{22. Apa itu Temperature dalam sampling LLM?}
\textbf{Jawaban:}
Parameter yang mengontrol keacakan output model.
\begin{itemize}
    \item \textbf{Low Temperature (misal 0.2):} Output lebih deterministik, fokus, dan konservatif (memilih probabilitas tertinggi).
    \item \textbf{High Temperature (misal 0.8):} Output lebih kreatif, beragam, dan berisiko (memilih probabilitas lebih rendah).
\end{itemize}

\subsection{23. Jelaskan konsep RAG (Retrieval-Augmented Generation)!}
\textbf{Jawaban:}
Teknik menggabungkan LLM dengan basis pengetahuan eksternal.
1. \textbf{Retrieval:} Sistem mencari dokumen relevan dari database (biasanya vector DB) berdasarkan query user.
2. \textbf{Augmentation:} Dokumen tersebut dimasukkan ke dalam prompt sebagai konteks.
3. \textbf{Generation:} LLM menjawab pertanyaan berdasarkan konteks tersebut.
Ini mengurangi halusinasi dan memungkinkan model mengakses data terkini/privat.

\subsection{24. Apa itu Hallucination pada AI?}
\textbf{Jawaban:}
Kondisi di mana LLM menghasilkan informasi yang tampak meyakinkan dan koheren secara tata bahasa, tetapi faktualnya salah atau tidak berdasar. Terjadi karena model hanya memprediksi kata berikutnya secara statistik, bukan memverifikasi fakta.

\subsection{25. Bagaimana cara kerja RLHF (Reinforcement Learning from Human Feedback)?}
\textbf{Jawaban:}
Proses menyelaraskan (aligning) LLM dengan preferensi manusia.
1. \textbf{SFT (Supervised Fine-Tuning):} Melatih model dasar dengan data demonstrasi berkualitas.
2. \textbf{Reward Modeling:} Melatih model terpisah untuk memprediksi skor kualitas jawaban berdasarkan peringkat manusia (A lebih baik dari B).
3. \textbf{PPO (Proximal Policy Optimization):} Menggunakan Reward Model untuk mengoptimalkan kebijakan LLM agar menghasilkan jawaban bernilai tinggi.

\subsection{26. Apa itu LoRA (Low-Rank Adaptation)?}
\textbf{Jawaban:}
Teknik fine-tuning yang efisien (PEFT). Alih-alih memperbarui semua parameter model (yang sangat besar), LoRA menyuntikkan matriks rank rendah yang dapat dilatih ke dalam setiap layer transformer dan membekukan bobot asli. Ini mengurangi jumlah parameter yang harus dilatih hingga >99\%, memungkinkan fine-tuning LLM di GPU konsumen.

\subsection{27. Apa itu Context Window?}
\textbf{Jawaban:}
Batas maksimum token (input + output) yang dapat diproses model dalam satu kali interaksi. Jika melebihi batas ini, model akan "lupa" bagian awal percakapan. Model terbaru seperti Claude 3 memiliki context window hingga 200k token (setara buku tebal).

\subsection{28. Jelaskan Prompt Engineering!}
\textbf{Jawaban:}
Seni dan teknik merancang input (prompt) untuk mengarahkan model AI menghasilkan output yang diinginkan secara optimal. Teknik mencakup: zero-shot, few-shot (memberikan contoh), chain-of-thought (meminta langkah berpikir), dan role-playing.

\subsection{29. Apa itu Vector Database?}
\textbf{Jawaban:}
Database yang dirancang untuk menyimpan dan mencari vektor (embedding) dimensi tinggi. Memungkinkan pencarian semantik (mencari makna, bukan kata kunci persis). Contoh: Pinecone, Milvus, Weaviate, Chroma. Penting untuk aplikasi RAG.

\subsection{30. Apa perbedaan antara Fine-Tuning dan RAG?}
\textbf{Jawaban:}
\begin{itemize}
    \item \textbf{Fine-Tuning:} Mengubah bobot internal model untuk mengajarkan \textit{cara/gaya} baru atau domain knowledge spesifik yang statis. "Mengajarkan siswa materi baru."
    \item \textbf{RAG:} Memberikan konteks eksternal ke dalam prompt tanpa mengubah model. Bagus untuk data yang sering berubah atau fakta spesifik. "Memberikan siswa buku teks saat ujian (open-book)."
\end{itemize}

\section{System Design & MLOps}

\subsection{31. Bagaimana Anda menangani dataset yang tidak seimbang (Imbalanced Dataset)?}
\textbf{Jawaban:}
\begin{itemize}
    \item \textbf{Resampling:} Undersampling kelas mayoritas atau Oversampling kelas minoritas (misal: SMOTE).
    \item \textbf{Class Weights:} Memberikan bobot lebih besar pada error di kelas minoritas dalam loss function.
    \item \textbf{Metrik yang Tepat:} Gunakan F1-Score, Precision-Recall AUC, bukan Akurasi.
    \item \textbf{Ensemble Methods:} Menggunakan Random Forest atau XGBoost yang lebih tahan banting.
\end{itemize}

\subsection{32. Bagaimana cara men-deploy model ML ke produksi?}
\textbf{Jawaban:}
1. \textbf{Model Serialization:} Simpan model (pickle, onnx, saved\_model).
2. \textbf{API Wrapper:} Bungkus model dalam REST API (FastAPI/Flask).
3. \textbf{Containerization:} Buat Docker image agar lingkungan konsisten.
4. \textbf{Orchestration:} Deploy menggunakan Kubernetes atau layanan cloud (AWS SageMaker, Google Vertex AI).
5. \textbf{Monitoring:} Pantau drift, latency, dan error rate.

\subsection{33. Apa itu Model Drift (Data Drift vs Concept Drift)?}
\textbf{Jawaban:}
Penurunan kinerja model seiring waktu.
\begin{itemize}
    \item \textbf{Data Drift:} Distribusi data input berubah (misal: sensor rusak, tren pengguna berubah). $P(X)$ berubah.
    \item \textbf{Concept Drift:} Hubungan antara input dan output berubah (misal: definisi "spam" berubah seiring waktu). $P(Y|X)$ berubah.
\end{itemize}

\subsection{34. Bagaimana Anda memilih jumlah cluster (K) dalam K-Means?}
\textbf{Jawaban:}
\begin{itemize}
    \item \textbf{Elbow Method:} Plot inersia (sum of squared distances) vs K. Pilih titik di mana penurunan inersia mulai melambat (membentuk siku).
    \item \textbf{Silhouette Score:} Mengukur seberapa mirip objek dengan clusternya sendiri dibandingkan dengan cluster lain.
    \item \textbf{Domain Knowledge:} Berdasarkan kebutuhan bisnis (misal: kita butuh 3 segmen pelanggan: Low, Mid, High).
\end{itemize}

\subsection{35. Apa itu A/B Testing dalam konteks ML?}
\textbf{Jawaban:}
Eksperimen terkontrol di mana kita membandingkan dua varian model secara langsung di lingkungan produksi.
\begin{itemize}
    \item \textbf{Group A (Control):} Dilayani oleh model lama (atau tanpa model).
    \item \textbf{Group B (Treatment):} Dilayani oleh model baru.
\end{itemize}
Metrik bisnis (CTR, konversi, revenue) diukur untuk menentukan pemenang secara statistik.

\section{Python & Coding}

\subsection{36. Apa perbedaan List dan Tuple di Python?}
\textbf{Jawaban:}
\begin{itemize}
    \item \textbf{List:} Mutable (bisa diubah isinya), menggunakan kurung siku \texttt{[]}. Lebih lambat.
    \item \textbf{Tuple:} Immutable (tidak bisa diubah setelah dibuat), menggunakan kurung biasa \texttt{()}. Lebih cepat dan memory-efficient. Bisa dipakai sebagai key dictionary.
\end{itemize}

\subsection{37. Jelaskan konsep Decorator di Python!}
\textbf{Jawaban:}
Fungsi yang menerima fungsi lain sebagai input dan mengembalikan fungsi baru dengan fungsionalitas tambahan, tanpa mengubah kode fungsi asli. Menggunakan sintaks \texttt{@decorator\_name}. Sering dipakai untuk logging, timing, atau autentikasi.

\subsection{38. Bagaimana cara menangani Missing Values di Pandas?}
\textbf{Jawaban:}
\begin{itemize}
    \item \textbf{Drop:} \texttt{df.dropna()} menghapus baris/kolom.
    \item \textbf{Imputasi Statistik:} \texttt{df.fillna(df.mean())} atau median/mode.
    \item \textbf{Forward/Backward Fill:} \texttt{df.fillna(method='ffill')}.
    \item \textbf{Model-based:} Menggunakan KNN Imputer atau regresi untuk memprediksi nilai yang hilang.
\end{itemize}

\subsection{39. Apa itu Generator di Python?}
\textbf{Jawaban:}
Fungsi yang mengembalikan iterator (menggunakan keyword \texttt{yield}) yang menghasilkan item satu per satu (lazy evaluation). Ini sangat hemat memori dibandingkan membuat list penuh, terutama untuk dataset besar.

\subsection{40. Jelaskan perbedaan \texttt{copy()} dan \texttt{deepcopy()}!}
\textbf{Jawaban:}
\begin{itemize}
    \item \textbf{Shallow Copy:} Membuat objek baru, tapi elemen di dalamnya (jika berupa objek mutabel seperti list) masih mereferensi ke objek asli. Perubahan pada elemen nested akan mempengaruhi aslinya.
    \item \textbf{Deep Copy:} Menyalin objek dan seluruh elemen nested di dalamnya secara rekursif. Objek baru benar-benar independen.
\end{itemize}

\section{Behavioral \& Situational}

\subsection{41. Ceritakan proyek AI tersulit yang pernah Anda kerjakan!}
\textbf{Tips:} Gunakan metode STAR (Situation, Task, Action, Result). Fokus pada tantangan teknis (misal: data kotor, resource terbatas) dan bagaimana Anda mengatasinya.

\subsection{42. Bagaimana Anda menjelaskan konsep teknis (misal Neural Network) kepada klien non-teknis?}
\textbf{Tips:} Gunakan analogi. "Neural network itu seperti anak kecil yang belajar mengenali hewan. Awalnya dia salah menebak kucing sebagai anjing, tapi setelah dikoreksi ribuan kali (latihan), dia mulai mengenali ciri khas masing-masing (telinga, ekor) dan menjadi ahli."

\subsection{43. Apa yang Anda lakukan jika model Anda memiliki performa bagus di training tapi buruk di produksi?}
\textbf{Jawaban:}
Ini indikasi overfitting atau data mismatch. Saya akan memeriksa:
1. Apakah distribusi data produksi berbeda dengan training (drift)?
2. Apakah ada kebocoran data (data leakage) saat training?
3. Saya akan mencoba mengumpulkan data baru dari produksi dan melatih ulang (retrain).

\subsection{44. Bagaimana Anda tetap update dengan perkembangan AI yang sangat cepat?}
\textbf{Jawaban:}
Membaca paper di ArXiv, mengikuti konferensi (NeurIPS, ICML), membaca newsletter (The Batch, TLDR AI), mengikuti peneliti top di Twitter/LinkedIn, dan bereksperimen langsung dengan model baru (Hugging Face).

\subsection{45. Apa pertimbangan etis yang Anda pikirkan saat membangun sistem AI?}
\textbf{Jawaban:}
Bias data (apakah adil untuk semua demografi?), privasi pengguna (GDPR), transparansi (explainability), dan potensi penyalahgunaan (dual-use).

\section{Studi Kasus}

\subsection{46. Bagaimana Anda mendesain sistem rekomendasi untuk e-commerce baru (cold start)?}
\textbf{Jawaban:}
Karena belum ada data interaksi user-item:
1. Gunakan \textbf{Content-based filtering} berdasarkan metadata produk.
2. Tampilkan produk populer/trending (non-personalized).
3. Minta preferensi eksplisit saat user mendaftar (onboarding).
4. Gunakan data kontekstual (lokasi, waktu, device).

\subsection{47. Bagaimana mendeteksi penipuan kartu kredit (Fraud Detection)?}
\textbf{Jawaban:}
Ini masalah klasifikasi highly imbalanced dan real-time.
1. \textbf{Data:} Transaksi historis (jumlah, waktu, lokasi, merchant).
2. \textbf{Model:} Anomaly Detection (Isolation Forest, Autoencoder) atau Supervised Learning (XGBoost) dengan sampling/class weights.
3. \textbf{Fitur:} Kecepatan transaksi (jumlah transaksi 1 jam terakhir), jarak geografis antar transaksi.

\subsection{48. Perusahaan ingin membuat chatbot customer service. Apa strateginya?}
\textbf{Jawaban:}
1. \textbf{Fase 1 (Rule-based):} Menangani FAQ sederhana.
2. \textbf{Fase 2 (RAG):} Gunakan LLM + Knowledge Base dokumen internal perusahaan untuk menjawab pertanyaan spesifik.
3. \textbf{Evaluasi:} Gunakan metrik relevansi dan human feedback.
4. \textbf{Safety:} Pasang guardrails agar bot tidak bicara kasar atau kompetitor.

\subsection{49. Bagaimana memprediksi harga saham?}
\textbf{Jawaban:}
(Hati-hati menjawab ini karena sangat sulit/stokastik).
Gunakan Time Series Analysis (ARIMA, Prophet) atau Deep Learning (LSTM, Transformer). Fitur: harga historis, volume, sentimen berita (NLP), indikator makroekonomi. Namun, tekankan bahwa pasar sangat noisy dan efisien.

\subsection{50. Bagaimana cara menghitung jumlah orang dalam kerumunan dari video CCTV?}
\textbf{Jawaban:}
Gunakan Computer Vision.
1. \textbf{Deteksi Objek:} YOLO atau Faster R-CNN untuk mendeteksi kepala/tubuh manusia.
2. \textbf{Density Map:} Jika sangat padat, gunakan regresi kepadatan (CSRNet).
3. \textbf{Tracking:} Gunakan algoritma tracking (DeepSORT) untuk menghitung orang unik yang melintas garis virtual.
